% !TEX root = ../diff_equations.tex
% =================================================================
%  sections/lecture02.tex
%  Лекция 2. Полные дифференциалы. Системы. Теорема Пеано. Условие Липшица.
% =================================================================

\part*{Лекция 2}

\section{Полные дифференциалы. Системы. Теорема Пеано. Условие Липшица}

\subsection{Уравнения в полных дифференциалах}\label{sec:полные-дифференциалы}

Пусть $F = m(x,y)\,\dd x + n(x,y)\,\dd y$.

\begin{Definition}{}{точная-форма}
	Форма $F$ называется \textbf{точной}, если
	\[
	\exists\, u(x,y) \in C^2(G) \colon \qquad
	F = \pder{u}{x}\,\dd x + \pder{u}{y}\,\dd y.
	\]
\end{Definition}

\begin{Definition}{}{полные-дифференциалы}
	Уравнение $F = 0$ с точной формой $F$ называется \textbf{уравнением
	в полных дифференциалах}.
\end{Definition}

\begin{Theorem}{Точная форма даёт интеграл}{точная-интеграл}
	Пусть $F$ --- точная форма. Тогда в окрестности любой точки
	$(x_0,y_0) \in G$ функция $u$ --- интеграл уравнения
	\[
	\dder{y}{x} = -\frac{m}{n}
	\qquad\text{или}\qquad
	\dder{x}{y} = -\frac{n}{m}.
	\]
\end{Theorem}

\begin{proof}
	Пусть $n(x_0,y_0) \neq 0$; тогда $n(x,y) \neq 0$ в некоторой окрестности
	$V$ точки $(x_0,y_0)$.

	Пусть $y(x)$ --- решение уравнения $y' = -\dfrac{m}{n}$ на
	$(\alpha,\beta)$, для которого $(x,y(x)) \in V$ при $x \in (\alpha,\beta)$.
	Тогда
	\[
	\dder{}{x} u\bigl(x,y(x)\bigr)
	= \pder{u}{x}\bigl(x,y(x)\bigr) + \pder{u}{y}\bigl(x,y(x)\bigr)\, y'(x)
	= m\bigl(x,y(x)\bigr) + n\bigl(x,y(x)\bigr)\cdot
	\Bigl(-\frac{m}{n}\bigl(x,y(x)\bigr)\Bigr) \equiv 0.
	\]
\end{proof}

\subsubsection*{Условие точности формы}

\begin{Statement}{Необходимое условие точности}{условие-точности}
	Если форма $F$ точна, то $m,n \in C^1$ и
	\begin{equation}\label{eq:условие-точности}
		\pder{m}{y}
		= \frac{\partial^2 u}{\partial x\,\partial y}
		= \frac{\partial^2 u}{\partial y\,\partial x}
		= \pder{n}{x}.
	\end{equation}
\end{Statement}

\begin{Remark}{Когда условие достаточно}{}
	Условие~\eqref{eq:условие-точности} достаточно, если $G$ --- односвязная
	область.
\end{Remark}

\subsubsection*{Частный случай: прямоугольник}

Пусть $G = (a,b)\times(\alpha,\beta)$ и $(x_0,y_0) \in G$
(рис.~\ref{fig:путь-интегрирования}).

Ищем $u$ из условия $\pder{u}{x} = m$:
\[
u = \int_{x_0}^{x} m(s,y)\,\dd s + \vphi(y), \qquad \vphi \in C^1.
\]

Тогда
\[
\pder{u}{y}
= \int_{x_0}^{x} \pder{m}{y}(s,y)\,\dd s + \vphi'(y)
= \int_{x_0}^{x} \pder{n}{x}(s,y)\,\dd s + \vphi'(y)
= n(x,y) - n(x_0,y) + \vphi'(y),
\]
где во втором равенстве использовано~\eqref{eq:условие-точности}.

Требуем, чтобы это выражение равнялось $n(x,y)$:
\[
n(x,y) - n(x_0,y) + \vphi'(y) = n(x,y)
\iff \vphi'(y) = n(x_0,y),
\qquad
\vphi(y) = \int_{y_0}^{y} n(x_0,t)\,\dd t.
\]

Итого
\[
u(x,y) = \int_{x_0}^{x} m(s,y)\,\dd s + \int_{y_0}^{y} n(x_0,t)\,\dd t.
\]

% ---------------------------------------------------------------
% РИСУНОК 1: прямоугольник G = (a,b) x (alpha,beta) и путь интегрирования:
% ломаная (x_0,y_0) -> (x_0,y) -> (x,y). Заливка области — hlBlue.
% ---------------------------------------------------------------
\begin{figure}[H]
	\centering
	\begin{tikzpicture}[x=1.4cm, y=1.3cm, >=Stealth]
		\draw[ax] (-0.35,0) -- (4.3,0) node[below right] {$x$};
		\draw[ax] (0,-0.35) -- (0,3.2) node[above] {$y$};

		% область G
		\fill[hlBlue, hlarea] (0.7,0.6) rectangle (3.6,2.7);
		\draw[curvebd] (0.7,0.6) rectangle (3.6,2.7);
		\node[slbl, below right] at (3.15,1.02) {$G$};

		% путь интегрирования: вверх по x = x_0, затем вправо
		\draw[curvehl, -{Stealth[length=2.4mm]}] (1.2,1.0) -- (1.2,2.1);
		\draw[curvehl, -{Stealth[length=2.4mm]}] (1.2,2.1) -- (2.9,2.1);

		% проекции на оси
		\draw[curveaux] (1.2,1.0) -- (0,1.0);
		\draw[curveaux] (1.2,2.1) -- (0,2.1);
		\draw[curveaux] (1.2,1.0) -- (1.2,0);
		\draw[curveaux] (2.9,2.1) -- (2.9,0);
		\node[slbl, left]  at (-0.06,1.0) {$y_0$};
		\node[slbl, left]  at (-0.06,2.1) {$y$};
		\node[slbl, below] at (1.2,-0.06) {$x_0$};
		\node[slbl, below] at (2.9,-0.06) {$x$};

		\node[pt] at (1.2,1.0) {};
		\node[pt] at (1.2,2.1) {};
		\node[pt] at (2.9,2.1) {};
		\node[slbl, below right] at (1.16,1.05) {$(x_0,y_0)$};
		\node[slbl, above right] at (1.16,2.12) {$(x_0,y)$};
		\node[slbl, above right] at (2.7,2.12) {$(x,y)$};
	\end{tikzpicture}
	\caption{Прямоугольная область $G = (a,b)\times(\alpha,\beta)$ и путь
		интегрирования: сначала вдоль прямой $x = x_0$ от $y_0$ до $y$, затем
		по горизонтали до точки $(x,y)$. Первому участку отвечает слагаемое
		$\int_{y_0}^{y} n(x_0,t)\,\dd t$, второму ---
		$\int_{x_0}^{x} m(s,y)\,\dd s$.}
	\label{fig:путь-интегрирования}
\end{figure}

\begin{Example}[addon]{Как работает формула}{полные-дифференциалы-пример}
	Пусть
	\[
	2xy\,\dd x + (x^2 + 3y^2)\,\dd y = 0 .
	\]
	Здесь $m = 2xy$, $n = x^2 + 3y^2$, и условие точности выполнено:
	$\pder{m}{y} = 2x = \pder{n}{x}$. Возьмём $(x_0,y_0) = (0,0)$; формула даёт
	\[
	u(x,y) = \int_0^x 2sy\,\dd s + \int_0^y 3t^2\,\dd t = x^2 y + y^3 .
	\]

	\medskip
	Проверка: $\pder{u}{x} = 2xy = m$ и $\pder{u}{y} = x^2 + 3y^2 = n$. Общий
	интеграл --- семейство кривых $x^2 y + y^3 = C$; условие
	$\pder{u}{y} \neq 0$ выполнено всюду, кроме начала координат.
\end{Example}

\begin{Example}[addon]{Замкнутая форма, не являющаяся точной}{неточная-замкнутая}
	Условие~\eqref{eq:условие-точности} не достаточно, если область не
	односвязна. Возьмём
	\[
	F = \frac{-y\,\dd x + x\,\dd y}{x^2 + y^2}, \qquad
	G = \RR^2 \setminus \{(0,0)\}.
	\]
	Здесь $m = -\dfrac{y}{x^2+y^2}$, $n = \dfrac{x}{x^2+y^2}$, и во всех
	точках $G$
	\[
	\pder{m}{y} = \frac{y^2 - x^2}{(x^2+y^2)^2} = \pder{n}{x}.
	\]

	\medskip
	Тем не менее функции $u \in C^2(G)$ с
	$F = \pder{u}{x}\,\dd x + \pder{u}{y}\,\dd y$ не существует. Возьмём
	окружность $\gamma(t) = (R\cos t,\, R\sin t)$, $t \in [0,\ 2\pi]$. Прямой
	подсчёт даёт
	\[
	m\bigl(\gamma(t)\bigr)\dot\gamma_1(t)
	+ n\bigl(\gamma(t)\bigr)\dot\gamma_2(t)
	= \frac{R^2\sin^2 t + R^2\cos^2 t}{R^2} \equiv 1.
	\]
	Если бы $F$ была точна, левая часть равнялась бы
	$\dder{}{t} u\bigl(\gamma(t)\bigr)$, и интегрирование по $[0,\ 2\pi]$
	дало бы одновременно
	\[
	\int_0^{2\pi} 1\,\dd t = 2\pi
	\qquad\text{и}\qquad
	u\bigl(\gamma(2\pi)\bigr) - u\bigl(\gamma(0)\bigr) = 0
	\]
	--- противоречие.

	\medskip
	Локально первообразная есть: в полуплоскости $x > 0$ подходит
	$u = \arctan\dfrac{y}{x}$, то есть полярный угол. Но продолжить её на
	всю область однозначно нельзя: при обходе вокруг нуля угол прирастает
	на $2\pi$.
\end{Example}

\subsubsection*{Интегрирующий множитель}

\begin{Definition}{}{интегрирующий-множитель}
	Функция $\mu(x,y) \in C^1$, $\mu \neq 0$, называется
	\textbf{интегрирующим множителем} формы $F = m\,\dd x + n\,\dd y$, если
	форма $\mu F = \mu m\,\dd x + \mu n\,\dd y$ точна. По
	условию~\eqref{eq:условие-точности} это равносильно тождеству
	\[
	\pder{}{y}(\mu m) \equiv \pder{}{x}(\mu n).
	\]
\end{Definition}

\begin{Example}[addon]{Два интегрирующих множителя}{интегрирующий-множитель-примеры}
	\textbf{1.} Форма $F = y\,\dd x - x\,\dd y$ не точна:
	$\pder{m}{y} = 1$, а $\pder{n}{x} = -1$. Где $y \neq 0$, возьмём
	$\mu = 1/y^2$:
	\[
	\mu F = \frac{\dd x}{y} - \frac{x\,\dd y}{y^2} = \dd\Bigl(\frac{x}{y}\Bigr),
	\]
	так что $u = x/y$ и общий интеграл $x/y = C$ --- прямые, проходящие через
	начало координат.

	\medskip
	\textbf{2.} Линейное уравнение $y' = p(x)y + q(x)$ в симметричной форме --- это
	$\bigl(p(x)y + q(x)\bigr)\dd x - \dd y = 0$, и оно не точно при $p \neq 0$.
	Множитель $\mu = e^{-\int p(x)\,\dd x}$ зависит только от $x$ и подходит:
	\[
	\pder{}{y}(\mu m) = p\,e^{-\int p\,\dd x}
	= \pder{}{x}\bigl(-e^{-\int p\,\dd x}\bigr) = \pder{}{x}(\mu n).
	\]
	Восстанавливая $u$, получаем
	$u = -y\,e^{-\int p\,\dd x} + \int q\,e^{-\int p\,\dd x}\,\dd x$, и равенство
	$u = \const$ --- это в точности формула Лагранжа из лекции~1. Метод вариации
	постоянной оказывается частным случаем интегрирующего множителя.
\end{Example}

\begin{Remark}[addon]{Множитель не единственен}{множитель-не-единственен}
	Если $\mu$ --- интегрирующий множитель и $\mu F = \dd u$, то для любой
	$\Phi \in C^1$, не обращающейся в нуль, форма
	$\Phi(u)\,\mu F = \Phi(u)\,\dd u$ тоже точна: это дифференциал
	первообразной функции $\Phi$. Значит, множителей всегда бесконечно много.

	\medskip
	Локально интегрирующий множитель существует всегда, но пользы от этого мало:
	его отыскание в общем случае не проще, чем решение самого уравнения.
\end{Remark}

\subsection{Системы дифференциальных уравнений}\label{sec:системы}

Пусть $x_1(t), \dots, x_n(t)$ --- искомые функции, а $m_1, \dots, m_n$ ---
натуральные числа.

\begin{Definition}{}{система-старшие}
	\textbf{Системой, разрешённой относительно старших производных},
	называется система
	\[
	x_i^{(m_i)} = f_i\bigl(t,\, x_1, \dot x_1, \dots, x_1^{(m_1-1)},\,
	\dots,\, x_n, \dot x_n, \dots, x_n^{(m_n-1)}\bigr),
	\qquad i = 1,\dots,n.
	\]
	Число $m = m_1 + \dots + m_n$ называется \textbf{порядком} системы.
\end{Definition}

\subsubsection*{Частные случаи}

\begin{Definition}{}{нормальная-система}
	При $m_1 = \dots = m_n = 1$ система называется \textbf{нормальной}:
	\[
	\dot x_i = f_i(t, x_1, \dots, x_n), \qquad i = 1,\dots,n.
	\]
\end{Definition}

При $n = 1$ получается дифференциальное уравнение порядка $m$:
\begin{equation}\label{eq:уравнение-порядка-m}
	x^{(m)} = f\bigl(t,\, x, \dot x, \dots, x^{(m-1)}\bigr).
\end{equation}

\subsubsection*{Сведение уравнения порядка $m$ к нормальной системе}

Вспомогательная система с искомыми функциями $y_1, \dots, y_m$:
\begin{equation}\label{eq:вспомогательная-система}
	\begin{cases}
		\dot y_1 = y_2,\\
		\dot y_2 = y_3,\\
		\ \ \vdots\\
		\dot y_m = f(t, y_1, \dots, y_m).
	\end{cases}
\end{equation}

\begin{Statement}{Соответствие решений}{сведение-к-системе}
	\begin{enumerate}
		\item Если $x(t)$ --- решение~\eqref{eq:уравнение-порядка-m}, то
		$y_1 = x$, $y_2 = \dot x$, \dots, $y_m = x^{(m-1)}$ ---
		решение~\eqref{eq:вспомогательная-система}.
		\item Если $y_1(t), \dots, y_m(t)$ ---
		решение~\eqref{eq:вспомогательная-система}, то $x = y_1$ ---
		решение~\eqref{eq:уравнение-порядка-m}, причём $\dot x = y_2$, \dots,
		$x^{(m-1)} = y_m$.
	\end{enumerate}
\end{Statement}

\begin{Example}[addon]{Гармонический осциллятор}{сведение-пример}
	Уравнение $\ddot x + \omega^2 x = 0$ имеет порядок $m = 2$. Полагая
	$y_1 = x$, $y_2 = \dot x$, получаем нормальную систему
	\[
	\begin{cases}
		\dot y_1 = y_2,\\
		\dot y_2 = -\omega^2 y_1 .
	\end{cases}
	\]
	По утверждению~\ref{stmt:сведение-к-системе} её решения и решения исходного
	уравнения --- одно и то же: $x = y_1 = C_1\cos\omega t + C_2\sin\omega t$, а
	$y_2 = \dot x$ восстанавливается дифференцированием.
\end{Example}

\subsection{Векторная запись}\label{sec:векторная-запись}

Дальше рассматриваются нормальные системы. Обозначения:
\[
x = \begin{pmatrix} x_1 \\ \vdots \\ x_n \end{pmatrix} \in \RR^n,
\qquad
\dot x = \begin{pmatrix} \dot x_1 \\ \vdots \\ \dot x_n \end{pmatrix},
\qquad
\int x(t)\,\dd t
= \begin{pmatrix} \int x_1\,\dd t \\ \vdots \\ \int x_n\,\dd t \end{pmatrix}.
\]

Будем рассматривать норму $\abs{x} = \max_{1 \leqslant i \leqslant n} \abs{x_i}.$

\begin{Remark}[addon]{Выбор нормы ни на что не влияет}{выбор-нормы}
	Все нормы в $\RR^n$ эквивалентны: для любых двух найдутся $c, C > 0$ с
	$c\norma{x}_1 \leqslant \norma{x}_2 \leqslant C\norma{x}_1$. Поэтому
	ограниченность, сходимость, непрерывность и условие Липшица не зависят от
	того, какая норма выбрана, --- меняются только константы в оценках.
	Максимум-норма берётся ради удобства: покоординатные оценки в ней
	автоматически дают оценку вектора.
\end{Remark}

\begin{Statement}{Оценка интеграла}{оценка-интеграла}
	\[
	\abs{\int_a^b x(t)\,\dd t} \leqslant \abs{\int_a^b \abs{x(t)}\,\dd t}.
	\]
\end{Statement}

\begin{Remark}[addon]{Зачем модуль справа}{внешний-модуль}
	Здесь не предполагается $a \leqslant b$, а при $b < a$ интеграл от
	неотрицательной функции $\abs{x(t)}$ отрицателен --- поэтому справа и стоит
	внешний модуль.
\end{Remark}

Полагая
\[
f(t,x) = \begin{pmatrix}
	f_1(t, x_1, \dots, x_n) \\ \vdots \\ f_n(t, x_1, \dots, x_n)
\end{pmatrix},
\]
нормальную систему записывают в векторном виде
\begin{equation}\label{eq:нормальная-система}
	\dot x = f(t,x), \qquad f \in C(G), \qquad G \subset \RR^{n+1}_{t,x}.
\end{equation}

\begin{Definition}{}{решение-системы}
	Функция $x \colon (a,b) \to \RR^n$ называется \textbf{решением}
	системы~\eqref{eq:нормальная-система}, если
	\begin{enumerate}
		\item $\exists\, \dot x(t)$, $t \in (a,b)$;
		\item $(t, x(t)) \in G$, $t \in (a,b)$;
		\item $\dot x(t) \equiv f\bigl(t, x(t)\bigr)$, $t \in (a,b)$.
	\end{enumerate}
\end{Definition}

\begin{Definition}{}{коши-система}
	Пусть $(t_0, x_0) \in G$. Функция $x(t)$ --- \textbf{решение задачи Коши}
	с начальными данными $(t_0,x_0)$, если
	\begin{enumerate}
		\item $x(t)$ --- решение на некотором $(a,b) \ni t_0$;
		\item $x(t_0) = x_0$.
	\end{enumerate}
\end{Definition}

\begin{Definition}{}{точка-единственности-система}
	Точка $(t_0,x_0) \in G$ --- \textbf{точка единственности}, если для любых
	решений $x_1(t)$, $x_2(t)$ задачи Коши с начальными данными $(t_0,x_0)$
	\[
	\exists\,(\alpha,\beta) \ni t_0 \colon \quad
	x_1(t) \equiv x_2(t), \quad t \in (\alpha,\beta).
	\]
\end{Definition}

\subsection{Теорема Пеано}\label{sec:пеано}

\begin{Theorem}{Пеано}{пеано}
	Пусть $f \in C(G)$. Тогда для любой точки $(t_0,x_0) \in G$ существует
	решение задачи Коши с начальными данными $(t_0,x_0)$.
\end{Theorem}

\subsubsection*{Эквивалентное интегральное уравнение}

\begin{equation}\label{eq:эиу}
	x(t) = x_0 + \int_{t_0}^{t} f\bigl(s, x(s)\bigr)\,\dd s.
\end{equation}

\begin{Definition}{}{решение-эиу}
	Функция $x \colon (a,b) \to \RR^n$ --- \textbf{решение}
	уравнения~\eqref{eq:эиу}, если
	\begin{enumerate}
		\item $x \in C(a,b)$;
		\item $(t, x(t)) \in G$, $t \in (a,b)$;
		\item равенство~\eqref{eq:эиу} выполнено на $(a,b)$.
	\end{enumerate}
\end{Definition}

\begin{Lemma}{Эквивалентность задачи Коши и интегрального уравнения}{эквивалентность-эиу}
	Функция $x(t)$ --- решение задачи Коши с начальными данными $(t_0,x_0)$
	на $(a,b)$ тогда и только тогда, когда $x(t)$ --- решение
	уравнения~\eqref{eq:эиу} на $(a,b)$.
\end{Lemma}

\begin{proof}
	\textbf{Необходимость.} Пункты 1--2 определения решения задачи Коши дают
	пункты 1--2 определения решения~\eqref{eq:эиу}. Далее,
	$\dot x(t) = f\bigl(t, x(t)\bigr)$ при $t \in (a,b)$; интегрируя от $t_0$
	до $t$, получаем
	\[
	\int_{t_0}^{t} \dot x(s)\,\dd s = \int_{t_0}^{t} f\bigl(s, x(s)\bigr)\,\dd s,
	\qquad
	x(t) - x(t_0) = \int_{t_0}^{t} f\bigl(s, x(s)\bigr)\,\dd s,
	\]
	а так как $x(t_0) = x_0$, это и есть~\eqref{eq:эиу}.

	\medskip
	\textbf{Достаточность.} Из $x \in C(a,b)$ следует, что
	$f\bigl(t, x(t)\bigr)$ непрерывна на $(a,b)$, поэтому интеграл
	$\int_{t_0}^{t} f\bigl(s,x(s)\bigr)\,\dd s$ принадлежит $C^1(a,b)$ ---
	это пункт 1 определения решения. Пункт 2 совпадает с пунктом 2
	определения решения~\eqref{eq:эиу}. Дифференцируя~\eqref{eq:эиу},
	получаем $\dot x(t) = f\bigl(t, x(t)\bigr)$, а подстановка $t = t_0$ даёт
	$x(t_0) = x_0$.
\end{proof}

\subsubsection*{Промежуток Пеано}

\begin{Definition}{}{промежуток-пеано}
	Пусть $\alpha, \beta > 0$ таковы, что
	\[
	R = \bigl\{(t,x) \colon \abs{t - t_0} \leqslant \alpha,\
	\abs{x - x_0} \leqslant \beta\bigr\} \subset G,
	\]
	и пусть $M > 0$ таково, что $\abs{f(t,x)} \leqslant M$ в $R$. Положим
	\[
	h = \min\Bigl(\alpha,\ \frac{\beta}{M}\Bigr).
	\]
	Отрезок $[t_0 - h,\ t_0 + h]$ называется \textbf{промежутком Пеано}
	(рис.~\ref{fig:промежуток-пеано}).
\end{Definition}

% ---------------------------------------------------------------
% РИСУНОК 2: прямоугольник R вокруг (t_0,x_0), конус наклонов
% |x-x_0| <= M|t-t_0| и промежуток Пеано на оси t.
% Параметры: t_0 = 2.4, x_0 = 1.6, alpha = 1.5, beta = 0.9, M = 1,
% так что h = beta/M = 0.9 < alpha. Масштабы по осям равны: важен наклон M.
% ---------------------------------------------------------------
\begin{figure}[H]
	\centering
	\begin{tikzpicture}[x=1.6cm, y=1.6cm, >=Stealth]
		\draw[ax] (-0.3,0) -- (4.6,0) node[below right] {$t$};
		\draw[ax] (0,-0.3) -- (0,3.05) node[above] {$x$};

		% прямоугольник R
		\fill[hlBlue, hlarea] (0.9,0.7) rectangle (3.9,2.5);

		% конус |x - x_0| <= M|t - t_0|, пересечённый с R
		\fill[hlPurple, hlarea] (0.9,0.7) -- (1.5,0.7) -- (2.4,1.6)
			-- (1.5,2.5) -- (0.9,2.5) -- cycle;
		\fill[hlPurple, hlarea] (3.9,0.7) -- (3.3,0.7) -- (2.4,1.6)
			-- (3.3,2.5) -- (3.9,2.5) -- cycle;

		% R замкнут, поэтому граница сплошная
		\draw[region] (0.9,0.7) rectangle (3.9,2.5);
		\node[slbl, below left] at (3.85,2.45) {$R$};

		% выносные линии
		\draw[curveaux] (0,1.6) -- (3.9,1.6);
		\foreach \t in {1.5,2.4,3.3} \draw[curveaux] (\t,0) -- (\t,2.5);
		\draw[curveaux] (0.9,0) -- (0.9,0.7);
		\draw[curveaux] (3.9,0) -- (3.9,0.7);
		\draw[curveaux] (0,0.7) -- (0.9,0.7);
		\draw[curveaux] (0,2.5) -- (0.9,2.5);

		% промежуток Пеано
		\draw[curvehl] (1.5,0) -- (3.3,0);

		\node[pt] at (2.4,1.6) {};
		\node[slbl, above right] at (2.46,1.66) {$(t_0,x_0)$};

		\foreach \y in {0.7,1.6,2.5} \draw[tickmark] (0.07,\y) -- (-0.07,\y);
		\node[slbl, left] at (-0.09,2.5) {$x_0+\beta$};
		\node[slbl, left] at (-0.09,1.6) {$x_0$};
		\node[slbl, left] at (-0.09,0.7) {$x_0-\beta$};

		\node[slbl, above, fill=white, inner sep=1pt] at (1.5,0.12) {$t_0-h$};
		\node[slbl, above, fill=white, inner sep=1pt] at (2.4,0.12) {$t_0$};
		\node[slbl, above, fill=white, inner sep=1pt] at (3.3,0.12) {$t_0+h$};
		\node[slbl, below] at (0.9,-0.12) {$t_0-\alpha$};
		\node[slbl, below] at (3.9,-0.12) {$t_0+\alpha$};
	\end{tikzpicture}
	\caption{Прямоугольник $R$ (голубой) и конус
		$\abs{x - x_0} \leqslant M\abs{t - t_0}$ (фиолетовый): график решения
		не может выйти из конуса, потому что $\abs{\dot x} \leqslant M$. Конус
		покидает $R$ через верхнюю и нижнюю стороны при
		$\abs{t - t_0} = \beta/M$, поэтому длина промежутка Пеано
		(выделенный отрезок на оси $t$) равна $h = \min(\alpha, \beta/M)$.}
	\label{fig:промежуток-пеано}
\end{figure}

\begin{Example}[addon]{Почему промежуток конечен}{промежуток-конечен}
	Пусть $\dot x = 1 + x^2$, $x(0) = 0$. Правая часть бесконечно гладкая на
	всей плоскости, а решение
	\[
	x(t) = \tan t
	\]
	определено лишь при $\abs{t} < \pi/2$. Решение задачи Коши существует только локально,
	и теорема Пеано большего не утверждает.

	\medskip
	Посмотрим, что даёт само построение. В прямоугольнике
	$R = \bigl\{\abs{t} \leqslant \alpha,\ \abs{x} \leqslant \beta\bigr\}$
	имеем $M = 1 + \beta^2$, поэтому
	\[
	h = \min\Bigl(\alpha,\ \frac{\beta}{1+\beta^2}\Bigr)
	\leqslant \max_{\beta > 0} \frac{\beta}{1+\beta^2} = \frac12 ,
	\]
	то есть никакой выбор $\alpha$ и $\beta$ не даёт $h > 1/2$. То есть оценка Пеано
	здесь грубее истинной ширины промежутка $\pi/2$.
\end{Example}

% ---------------------------------------------------------------
% РИСУНОК 3: решение x = tan t задачи Коши для x' = 1 + x^2, x(0) = 0,
% и вертикальные асимптоты t = ± pi/2 — границы его существования.
% ---------------------------------------------------------------
\begin{figure}[H]
	\centering
	\begin{tikzpicture}[x=1.55cm, y=0.6cm, >=Stealth]
		\draw[ax] (-2.15,0) -- (2.3,0) node[below right] {$t$};
		\draw[ax] (0,-3.5) -- (0,3.7) node[above] {$x$};

		% асимптоты t = ± pi/2
		\draw[curveaux] (-1.5708,-3.4) -- (-1.5708,3.5);
		\draw[curveaux] (1.5708,-3.4) -- (1.5708,3.5);

		% решение x = tan t
		\begin{scope}
			\clip (-1.8,-3.4) rectangle (1.8,3.4);
			\draw[curvehl, domain=-1.48:1.48, samples=200]
			plot (\x,{tan(\x r)});
		\end{scope}

		\node[pt] at (0,0) {};
		\foreach \t in {-1.5708,1.5708} \draw[tickmark] (\t,0.14) -- (\t,-0.14);
		\node[slbl, below left]  at (-1.52,-0.2) {$-\pi/2$};
		\node[slbl, below right] at (1.62,-0.2) {$\pi/2$};
	\end{tikzpicture}
	\caption{Решение задачи Коши $\dot x = 1 + x^2$, $x(0) = 0$: правая часть
		гладка во всей плоскости, а решение $x = \tan t$ живёт лишь между
		асимптотами $t = \pm\pi/2$ (пунктир). Продолжить его нельзя --- поэтому
		в теореме Пеано и появляется конечный промежуток.}
	\label{fig:тангенс}
\end{figure}

\begin{Statement}{Существование решения интегрального уравнения}{существование-эиу}
	Существует решение уравнения~\eqref{eq:эиу} на промежутке
	$[t_0,\ t_0 + h]$.
\end{Statement}

\subsubsection*{Ломаные Эйлера}

Фиксируем натуральное $N$ и положим
\[
t_k = t_0 + \frac{k}{N}\,h, \qquad k = 0, \dots, N.
\]

\textbf{Ломаная Эйлера} $g(t) = g_N(t)$ строится по звеньям
(рис.~\ref{fig:ломаная-эйлера}). При $t \in [t_0, t_1]$
\[
g(t) = x_0 + f(t_0, x_0)(t - t_0),
\]
а если $k < N$ и $(t_k, g(t_k)) \in G$, то
\[
g(t) = g(t_k) + f\bigl(t_k, g(t_k)\bigr)(t - t_k),
\qquad t \in [t_k, t_{k+1}].
\]

Производная в узлах определяется равенствами
\[
\dot g(t_0) = f(t_0, x_0),
\qquad
\dot g(t_k) = f\bigl(t_k, g(t_k)\bigr)
\quad\text{при } k > 0,
\]
если $g$ определена на $[t_k, t_{k+1}]$.

% ---------------------------------------------------------------
% РИСУНОК 4: ломаная Эйлера: три звена с разными наклонами,
% узлы t_0 < t_1 < t_2 < t_3 равноотстоят с шагом h/N.
% ---------------------------------------------------------------
\begin{figure}[H]
	\centering
	\begin{tikzpicture}[x=1.3cm, y=1.3cm, >=Stealth]
		\draw[ax] (-0.4,0) -- (4.4,0) node[below right] {$t$};
		\draw[ax] (0,-0.4) -- (0,2.6) node[above] {$x$};

		% проекции узлов на ось t
		\foreach \t/\lab in {0.6/{$t_0$}, 1.6/{$t_1$}, 2.6/{$t_2$}, 3.6/{$t_3$}}{
			\draw[curveaux] (\t,0) -- (\t,{0.8 + 0.45*(\t-0.6)});
			\node[slbl, below] at (\t,-0.06) {\lab};
		}
		\draw[curveaux] (0.6,0.8) -- (0,0.8);
		\node[slbl, left] at (-0.06,0.8) {$x_0$};

		% звенья ломаной
		\draw[curvehl] (0.6,0.8) -- (1.6,1.25) -- (2.6,1.55) -- (3.6,2.25);
		\foreach \p in {(0.6,0.8), (1.6,1.25), (2.6,1.55), (3.6,2.25)}{
			\node[pt] at \p {};
		}

		\node[slbl, above left] at (3.55,2.3) {$g_N$};
	\end{tikzpicture}
	\caption{Ломаная Эйлера: на каждом отрезке $[t_k, t_{k+1}]$ она идёт по
		прямой с угловым коэффициентом $f\bigl(t_k, g(t_k)\bigr)$, вычисленным
		в левом конце отрезка. Узлы делят промежуток Пеано на $N$ равных
		частей длины $h/N$.}
	\label{fig:ломаная-эйлера}
\end{figure}

\begin{Lemma}{Свойства ломаной Эйлера}{ломаная-свойства}
	Для всех $k = 1, \dots, N$:
	\begin{enumerate}
		\item $g(t)$ определена на $[t_0, t_k]$;
		\item $\abs{g(t) - x_0} \leqslant M\abs{t - t_0}$ на $[t_0, t_k]$;
		\item $g(t) = x_0 + \displaystyle\int_{t_0}^{t} \dot g(s)\,\dd s$
		на $[t_0, t_k]$.
	\end{enumerate}
\end{Lemma}

\begin{proof}
	Индукция по $k$; при $k = 1$ все три пункта очевидны. Пусть они верны
	на $[t_0, t_k]$; докажем их на $[t_0, t_{k+1}]$. Обозначим
	$g_k = g(t_k)$.

	\medskip
	\textbf{Пункт 1.} Так как $t_k - t_0 \leqslant h \leqslant \alpha$ и
	\[
	\abs{g_k - x_0} \leqslant M\abs{t_k - t_0} \leqslant Mh \leqslant \beta,
	\]
	точка $(t_k, g_k)$ лежит в $R \subset G$. Значит, $f(t_k, g_k)$
	определена, $\abs{f(t_k,g_k)} \leqslant M$, и $g(t)$ определена на
	$[t_k, t_{k+1}]$.

	\medskip
	\textbf{Пункт 2.} При $t \in [t_k, t_{k+1}]$
	\[
	\abs{g(t) - x_0} \leqslant M(t - t_k) + M(t_k - t_0) = M(t - t_0).
	\]

	\medskip
	\textbf{Пункт 3.} При $t \in [t_k, t_{k+1}]$
	\[
	g(t) = g(t_k) + f\bigl(t_k, g(t_k)\bigr)(t - t_k)
	= x_0 + \int_{t_0}^{t_k} \dot g(s)\,\dd s + \int_{t_k}^{t} \dot g(s)\,\dd s
	= x_0 + \int_{t_0}^{t} \dot g(s)\,\dd s.
	\]
\end{proof}

\subsubsection*{Лемма Арцела --- Асколи}

\begin{Definition}{}{арцела-условия}
	Последовательность $g_N(t)$, $N = 1, 2, \dots$, на отрезке $[a,b]$
	называется \textbf{равномерно ограниченной}, если
	\[
	\exists\, K \colon \quad \abs{g_N(t)} \leqslant K, \quad t \in [a,b],
	\quad N = 1, 2, \dots,
	\]
	и \textbf{равностепенно непрерывной}, если
	\[
	\forall \eps > 0\ \exists\, \delta > 0 \colon \quad
	t, t' \in [a,b],\ \abs{t - t'} \leqslant \delta
	\implies \abs{g_N(t) - g_N(t')} \leqslant \eps \quad \forall N.
	\]
\end{Definition}

\begin{Lemma}{Арцела --- Асколи}{арцела-асколи}
	Если последовательность $\{g_N(t)\}$ равномерно ограничена и
	равностепенно непрерывна на $[a,b]$, то из неё можно выбрать равномерно
	сходящуюся подпоследовательность.
\end{Lemma}

\begin{Lemma}{Ломаные Эйлера удовлетворяют условиям Арцела --- Асколи}{ломаные-компактность}
	Последовательность ломаных Эйлера равномерно ограничена и равностепенно
	непрерывна на $[t_0,\ t_0+h]$.
\end{Lemma}

\begin{proof}
	Для ломаных Эйлера $\abs{g_N(t) - x_0} \leqslant M(t - t_0) \leqslant Mh$, откуда
	\[
	\abs{g_N(t)} \leqslant \abs{x_0} + Mh,
	\]
	то есть последовательность равномерно ограничена.

	\medskip
	Пусть $\eps > 0$ и $\delta = \dfrac{\eps}{M}$. Тогда для
	$t, t' \in [t_0,\ t_0+h]$ с $\abs{t - t'} \leqslant \delta$
	\[
	\abs{g_N(t) - g_N(t')}
	= \abs{\int_{t}^{t'} \dot g_N(s)\,\dd s}
	\leqslant M\abs{t - t'} \leqslant M\delta = \eps.
	\]
\end{proof}

\subsubsection*{Предельная функция --- решение интегрального уравнения}

Считаем, что $g_N(t) \rightrightarrows g(t)$ на $[t_0,\ t_0+h]$
(рис.~\ref{fig:сходимость-ломаных}).

% ---------------------------------------------------------------
% РИСУНОК 5: ломаные Эйлера g_2, g_4, g_8 для x' = 1 + x^2, x(0) = 0
% на [0; 1,2] и их предел x = tan t. Узлы равномерны с шагом 1,2/N.
% ---------------------------------------------------------------
\begin{figure}[H]
	\centering
	\begin{tikzpicture}[x=3.7cm, y=1.25cm, >=Stealth]
		\draw[ax] (-0.1,0) -- (1.5,0) node[below right] {$t$};
		\draw[ax] (0,-0.15) -- (0,2.95) node[above] {$x$};

		% предельное решение x = tan t
		\draw[curvehl, domain=0:1.2, samples=120] plot (\x,{tan(\x r)});

		% ломаные Эйлера
		\draw[curve] (0,0) -- (0.6,0.6) -- (1.2,1.416);
		\draw[curve] (0,0) -- (0.3,0.3) -- (0.6,0.627) -- (0.9,1.045)
		-- (1.2,1.673);
		\draw[curve] (0,0) -- (0.15,0.15) -- (0.3,0.303) -- (0.45,0.467)
		-- (0.6,0.650) -- (0.75,0.863) -- (0.9,1.125) -- (1.05,1.465)
		-- (1.2,1.937);
		\foreach \q in {(0.6,0.6), (1.2,1.416)} \node[pt, inner sep=1pt] at \q {};
		\foreach \q in {(0.3,0.3), (0.6,0.627), (0.9,1.045), (1.2,1.673)} \node[pt, inner sep=1pt] at \q {};
		\foreach \q in {(0.15,0.15), (0.3,0.303), (0.45,0.467), (0.6,0.650), (0.75,0.863),
		(0.9,1.125), (1.05,1.465), (1.2,1.937)} \node[pt, inner sep=1pt] at \q {};

		\node[pt] at (0,0) {};
		\node[Purple, slbl, right] at (1.23,2.572) {$x = \tan t$};
		\node[Blue, slbl, right]   at (1.23,1.937) {$g_8$};
		\node[Blue, slbl, right]   at (1.23,1.673) {$g_4$};
		\node[Blue, slbl, right]   at (1.23,1.416) {$g_2$};
	\end{tikzpicture}
	\caption{Ломаные Эйлера $g_2$, $g_4$, $g_8$ для уравнения
		$\dot x = 1 + x^2$ с начальным условием $x(0) = 0$ при
		$0 \leqslant t \leqslant 1{,}2$. С измельчением разбиения они
		поднимаются к решению $x = \tan t$ (выделенная кривая). Лемма
		Арцела --- Асколи гарантирует лишь сходимость подпоследовательности;
		здесь начальная точка --- точка единственности, поэтому предел один и
		тот же для всей последовательности.}
	\label{fig:сходимость-ломаных}
\end{figure}

\begin{Statement}{Предел ломаных Эйлера}{предел-ломаных}
	Функция $g(t)$ --- решение уравнения~\eqref{eq:эиу} на $[t_0,\ t_0+h]$.
\end{Statement}

\textbf{Пункт 1.} Все $g_N \in C[t_0,\ t_0+h]$, поэтому и
$g \in C[t_0,\ t_0+h]$.

\textbf{Пункт 2.} Из $(t, g_N(t)) \in R$ следует $(t, g(t)) \in R \subset G$
при $t \in [t_0,\ t_0+h]$.

\textbf{Пункт 3.} Запишем
\[
g_N(t) = x_0 + \int_{t_0}^{t} \dot g_N(s)\,\dd s
= x_0
+ \underbrace{\int_{t_0}^{t} f\bigl(s, g_N(s)\bigr)\,\dd s}_{I_1}
+ \underbrace{\int_{t_0}^{t}
	\bigl[\dot g_N(s) - f\bigl(s, g_N(s)\bigr)\bigr]\,\dd s}_{I_2}.
\]

Так как $f\bigl(t, g_N(t)\bigr) \rightrightarrows f\bigl(t, g(t)\bigr)$ на
$[t_0,\ t_0+h]$, имеем $I_1 \to \int_{t_0}^{t} f\bigl(s, g(s)\bigr)\,\dd s$.
Остаётся проверить, что $I_2 \to 0$ при $N \to \infty$.

Фиксируем $\eps > 0$. Функция $f$ равномерно непрерывна на компакте $R$,
поэтому существует $\delta > 0$ такое, что
\[
(t,x), (t',x') \in R,\quad \abs{t - t'} \leqslant \delta,\
\abs{x - x'} \leqslant \delta
\implies \abs{f(t,x) - f(t',x')} \leqslant \eps.
\]

При больших $N$ выполнено $\dfrac{h}{N} < \delta$ и
$\dfrac{Mh}{N} < \delta$, поэтому для $s \in [t_k, t_{k+1})$
\[
\abs{s - t_k} < \delta,
\qquad
\abs{g_N(s) - g_N(t_k)} \leqslant M\abs{s - t_k} \leqslant \frac{M}{N}h < \delta.
\]

Но $\dot g_N(s) = f\bigl(t_k, g_N(t_k)\bigr)$ при $s \in [t_k, t_{k+1})$,
значит $\abs{\dot g_N(s) - f\bigl(s, g_N(s)\bigr)} < \eps$ и
$\abs{I_2} \leqslant \eps h \xrightarrow[N \to \infty]{} 0$.

Переходя в равенстве для $g_N(t)$ к пределу, получаем
\[
g(t) = x_0 + \int_{t_0}^{t} f\bigl(s, g(s)\bigr)\,\dd s,
\qquad t \in [t_0,\ t_0+h],
\]
то есть $g$ --- решение уравнения~\eqref{eq:эиу}; вместе с
леммой~\ref{lem:эквивалентность-эиу} это доказывает
теорему~\ref{thm:пеано}.

\begin{Remark}[addon]{Левая половина промежутка}{левая-половина}
	Построение проведено на правой половине $[t_0,\ t_0+h]$.
	Левая половина сводится к уже доказанному заменой
	$s = 2t_0 - t$: если положить $z(s) = x(2t_0 - s)$, то
	\[
	\dot z(s) = -\dot x(2t_0 - s) = -f\bigl(2t_0 - s,\ z(s)\bigr)
	=: \tilde f\bigl(s, z(s)\bigr),
	\]
	причём $\tilde f$ непрерывна и ограничена той же константой $M$ на
	прямоугольнике тех же размеров. Значит, к ней применимо доказанное, и
	решение $z$ на $[t_0,\ t_0+h]$ даёт решение $x(t) = z(2t_0 - t)$ на
	$[t_0-h,\ t_0]$.
\end{Remark}

\subsection{Условие Липшица}\label{sec:липшиц}

\begin{Definition}{}{липшиц}
	Говорят, что $f \in \mathrm{Lip}_x(G)$, если существует
	\textbf{константа Липшица} $L \geqslant 0$ такая, что
	\[
	\forall (t,x), (t,x') \in G \abs{f(t,x) - f(t,x')} \leqslant L\abs{x - x'}
	\]
	(рис.~\ref{fig:липшиц}).
\end{Definition}

% ---------------------------------------------------------------
% РИСУНОК 6: условие Липшица: две точки (t,x) и (t,x') на одной
% вертикали внутри области G.
% ---------------------------------------------------------------
\begin{figure}[H]
	\centering
	\begin{tikzpicture}[x=1.3cm, y=1.2cm, >=Stealth]
		\draw[ax] (-0.4,0) -- (4.3,0) node[below right] {$t$};
		\draw[ax] (0,-0.4) -- (0,3.1) node[above] {$x$};

		% область G
		\fill[hlBlue, hlarea] plot[smooth cycle, tension=0.85]
		coordinates {(0.8,0.8) (2.3,0.5) (3.7,1.1) (3.9,2.4) (2.4,3.0) (1.0,2.4)};
		\draw[curvebd] plot[smooth cycle, tension=0.85]
		coordinates {(0.8,0.8) (2.3,0.5) (3.7,1.1) (3.9,2.4) (2.4,3.0) (1.0,2.4)};
		\node[slbl] at (1.3,1.0) {$G$};

		\draw[curveaux] (2.3,2.25) -- (0,2.25);
		\draw[curveaux] (2.3,1.25) -- (0,1.25);
		\draw[curveaux] (2.3,2.6) -- (2.3,0);
		\node[slbl, left]  at (-0.06,2.25) {$x$};
		\node[slbl, left]  at (-0.06,1.25) {$x'$};
		\node[slbl, below] at (2.3,-0.06) {$t$};
		\node[slbl, right] at (2.35,1.75) {$\abs{x-x'}$};

		% вертикаль t = const и две точки на ней
		\draw[curvehl] (2.3,1.25) -- (2.3,2.25);
		\node[pt] at (2.3,2.25) {};
		\node[pt] at (2.3,1.25) {};
	\end{tikzpicture}
	\caption{Условие Липшица сравнивает значения $f$ в точках с одной и той же
		первой координатой: приращение $f$ по $x$ оценивается через расстояние
		$\abs{x - x'}$ с одной константой $L$ для всей области~$G$.}
	\label{fig:липшиц}
\end{figure}

\begin{Definition}{}{липшиц-лок}
	Говорят, что $f \in \mathrm{Lip}_{x,\mathrm{loc}}(G)$, если
	\[
	\forall (t_0,x_0) \in G\ \exists\ \text{окрестность } \mathcal{U}
	\text{ точки } (t_0,x_0) \colon \quad f \in \mathrm{Lip}_x(\mathcal{U}).
	\]
\end{Definition}

\begin{Definition}{}{операторная-норма}
	\textbf{Операторной нормой} матрицы $A$ называется
	\[
	\norma{A} = \max_{\abs{x} = 1} \abs{Ax},
	\qquad\text{откуда}\qquad
	\abs{Ax} \leqslant \norma{A}\cdot\abs{x} \quad \forall x.
	\]
	\textbf{Матрица Якоби} отображения $f$ по переменной $x$ ---
	\[
	\pder{f}{x} = \Bigl(\pder{f_i}{x_j}\Bigr), \qquad i,j = 1,\dots,n.
	\]
\end{Definition}

\begin{Remark}[addon]{Как её посчитать}{операторная-норма-формула}
	Определение через максимум по шару неконструктивно, но для выбранной
	максимум-нормы $\abs{x} = \max_i \abs{x_i}$ подчинённая операторная норма
	выписывается явно --- это наибольшая сумма модулей по строке:
	\[
	\norma{A} = \max_{1 \leqslant i \leqslant n} \sum_{j=1}^{n} \abs{a_{ij}} .
	\]
	Действительно, $\abs{(Ax)_i} \leqslant \sum_j \abs{a_{ij}}\abs{x_j}
	\leqslant \abs{x}\sum_j \abs{a_{ij}}$, а равенство достигается на векторе
	из $\pm 1$, знаки которого совпадают со знаками элементов нужной строки.
\end{Remark}

\begin{Lemma}{Непрерывная матрица Якоби даёт локальное условие Липшица}{якоби-липшиц}
	Если $\pder{f}{x} \in C(G)$, то $f \in \mathrm{Lip}_{x,\mathrm{loc}}(G)$.
\end{Lemma}

\begin{proof}
	Пусть $(t_0,x_0) \in G$ и $\alpha, \beta > 0$ таковы, что
	\[
	R = \bigl\{(t,x) \colon \abs{t - t_0} \leqslant \alpha,\
	\abs{x - x_0} \leqslant \beta\bigr\} \subset G.
	\]
	Из $\pder{f}{x} \in C(G)$ следует, что найдётся $L \geqslant 0$ с
	$\norma{\pder{f}{x}} \leqslant L$ в $R$.

	\medskip
	Положим $g(s) = f\bigl(t,\ sx + (1-s)x'\bigr)$, $s \in [0,1]$;
	прямоугольник $R$ выпукл, поэтому при $(t,x), (t,x') \in R$ весь отрезок
	$sx + (1-s)x'$ лежит в $R$ и $g$ определена. Тогда
	\[
	f(t,x) - f(t,x') = g(1) - g(0) = \int_0^1 \pder{g}{s}(s)\,\dd s
	= \int_0^1 \pder{f}{x}\bigl(t,\ sx + (1-s)x'\bigr)\cdot(x - x')\,\dd s,
	\]
	откуда
	\[
	\abs{f(t,x) - f(t,x')}
	\leqslant \abs{\int_0^1 \norma{\pder{f}{x}}\cdot\abs{x - x'}\,\dd s}
	\leqslant L\abs{x - x'}.
	\]
\end{proof}

\begin{Lemma}{Условие Липшица на компакте}{липшиц-компакт}
	Пусть $f \in C(G)$, $f \in \mathrm{Lip}_{x,\mathrm{loc}}(G)$ и
	$K \subset G$ --- компакт. Тогда $f \in \mathrm{Lip}_x(K)$.
\end{Lemma}

\begin{proof}
	Пусть это не так. Тогда для любого $L_n > 0$ найдутся точки
	$(t_n, x_n), (t_n, x'_n) \in K$, $n = 1, 2, \dots$, с
	\begin{equation}\label{eq:липшиц-противоречие}
		\abs{f(t_n, x_n) - f(t_n, x'_n)} > L_n \abs{x_n - x'_n}.
	\end{equation}
	Возьмём $L_n \to \infty$ при $n \to \infty$. Выберем сходящуюся
	подпоследовательность (пусть это сама $(t_n, x_n)$):
	$(t_n, x_n) \to (t_0, x_0) \in K$, а из неё --- подпоследовательность
	с $(t_n, x'_n) \to (t_0, x'_0) \in K$.

	\medskip
	\textbf{Случай 1: $x_0 = x'_0$.} У точки $(t_0, x_0) = (t_0, x'_0)$ есть
	окрестность $\mathcal{U}$, в которой $f \in \mathrm{Lip}_x(\mathcal{U})$,
	то есть существует $L$ с
	\[
	\abs{f(t,x) - f(t,x')} \leqslant L\abs{x - x'}
	\qquad\text{для } (t,x), (t,x') \in \mathcal{U}.
	\]
	При больших $n$ обе точки лежат в $\mathcal{U}$ --- противоречие
	с~\eqref{eq:липшиц-противоречие}.

	\medskip
	\textbf{Случай 2: $x_0 \neq x'_0$.} Рассмотрим
	\[
	g(t,x,y) = \frac{\abs{f(t,x) - f(t,y)}}{\abs{x - y}}.
	\]
	Функция $g$ непрерывна в окрестности точки $(t_0, x_0, x'_0)$, а потому
	ограничена в некоторой окрестности $\mathcal{U}$: существует $L > 0$ с
	\[
	\abs{f(t_n, x_n) - f(t_n, x'_n)} \leqslant L\abs{x_n - x'_n}
	\qquad\text{при больших } n,
	\]
	что снова противоречит~\eqref{eq:липшиц-противоречие}.
\end{proof}
